\documentclass[../main.tex]{subfiles}
\begin{document}
\begin{CJK*}{UTF8}{bkai}
\subsection{質量傳輸概論}
\begin{itemize}
  \item 何謂質量傳輸?\\
  質量傳輸就是有各種物質，在空間與時間下的分布與變化，因此必須要有
  \begin{itemize}
    \item 兩種以上的物料
    \item 物料的濃度梯度是傳輸的驅動力\\
    因此均勻濃度的流體並不會有質量傳輸
    \item 如果用擴散來表示質量傳輸時，與熱傳導概念相同
  \end{itemize}
  而擴散根據介質流動與否分成兩種:
  \begin{enumerate}
    \item molecular diffusion:(Steady-state diffusion)\\
    當介質\fbox{不流動}時，因物質濃度差異，而產生分子擴散\\
    固體的介質會慢於液體，液體又慢於氣體
    \begin{itemize}
      \item 以固體為介質的分子擴散，會用\fbox{Fick's Law}來描述
      \begin{equation}
        \boxed{J_{Az} = -D_{AB}\frac{dC_a}{dz}}
      \end{equation}
      e.g.
      \begin{itemize}
        \item $H_2$於$\text{Pd}_{s}$中擴散
        \item 尿素於洋菜中擴散
        \item 氦氣在玻璃中擴散
        \item 有機物於薄膜中擴散
      \end{itemize}
      \item 以液體或氣體為介質的分子擴散，會用\fbox{Fick's First Law}來描述
      \begin{equation}
        \boxed{N_{Az} = J_{Az}+ x_A(N_{Az}+N_{Bz}) = -C D_{AB} \frac{dx_A}{dz} + x_A (N_{Az}+N_{Bz})}
      \end{equation}
      \item 不管是什麼狀態，都可以使用的經驗式
      \begin{equation}
        \boxed{N_A = k_c\Delta C_A = k_y \Delta y_A = k_x \Delta x_A}
      \end{equation}
    \end{itemize}
    \item Mass convection:\\
    當介質\fbox{流動}時，物質會跟著介質一起流動，而產生質量傳送\\
    會類似於熱傳導中的對流
  \end{enumerate}
  \item 質量傳輸名詞定義\\
    假設空間中有一個控制體積，則在這個控制體積內\\
    各個物質的密度、速度、質量通量、以及生成速率，會隨著空間與時間的變化而改變
    \begin{enumerate}
      \item Density, $\rho_i(\vec{\bm r}, t)$，純量，對各方向皆相同\\
        若針對所有分密度相加，可以改為平均密度
        \begin{equation}
          \sum_i \rho_i = \rho
        \end{equation}
      \item Rate of Production, $r_i(\vec{\bm r}, t)$，純量，對各方向皆相同\\
        根據質量守恆，反應通量$r_i$，會滿足
        \begin{equation}
          \boxed{\sum_i r_i=0}
        \end{equation}
      \item Velocity, $\vec{\bm u}_i(\vec{\bm r}, t)$，向量，會隨方向改變
      \item Mass Flux, $\vec{\bm n}_i (\vec{\bm r}, t)$，向量，會隨方向改變\\
      質量通量，又可以再拆分為
      \begin{enumerate}
        \item 對流造成的質量通量\\
          在這個流場中，物質跟著整陀流體一起移動，而產生的質量通量
        \item 分子擴散造成的質量通量\\
          在這個流場中，物質因為濃度差異，額外產生的質量通量
      \end{enumerate}
      因此質量通量，可以表示為
      \begin{equation}
        \boxed{\vec{\bm n}_i = \rho_i \vec{\bm u} + \vec{\bm j}_i} \label{eq:mass_flux_decomposition} 
      \end{equation}
      其中\\
      $\vec{\bm u}$為\fbox{Volume Average Velocity}\\
      $\vec{\bm j}_i$即為分子擴散所造成的\fbox{Diffusive mass flux}
      \item Volume Average Velocity, $\vec{\bm u}$，向量，會隨方向改變\\
        針對所有物質的擴散加總，會得到一個淨流入與淨流出的速度，定義為體積平均流速
        \begin{equation}
          \sum_i \rho_i \vec{\bm u} = \rho \vec{\bm u} \implies \boxed{\vec{\bm u} = \frac{\sum_i \rho_i \vec{\bm u}}{\rho}}
        \end{equation}
      \item Diffusive mass flux, $\vec{\bm j}_i$，向量，會隨方向改變\\
        由(\ref{eq:mass_flux_decomposition})，各物質的擴散$\vec{\bm j}_i$\\
        則是根據平均流速為基準額外增加或減少的質量通量，使得整體
        \begin{equation}
          \boxed{\sum_i \vec{\bm j}_i = 0}
        \end{equation}
        推導如下:
        \begin{align}
          \vec{\bm j}_i &= \vec{\bm n}_i - \rho_i \vec{\bm u} \nonumber\\
          &= \rho_i (\vec{\bm u}_i - \vec{\bm u})\nonumber\\
          \sum_i \vec{\bm j}_i &= \sum_i \rho_i (\vec{\bm u}_i - \vec{\bm u})\nonumber\\
          &= \sum_i \rho_i \vec{\bm u}_i - \vec{\bm u} \sum_i \rho_i \nonumber\\
          &= \rho \vec{\bm u} - \rho \vec{\bm u} = 0
        \end{align}
      \item Molar 相關衍生單位，除上分子量後，以大寫表示
      \begin{enumerate}
        \item Species Molar Flux, $\vec{\bm N}_i = \frac{\vec{\bm n}_i}{M_i}$，向量，會隨方向改變\\
          對比Mass Flux (\ref{eq:mass_flux_decomposition})
          \begin{align}
            &\boxed{\frac{\vec{\bm n}_i}{M_i}} = \boxed{\frac{\rho_i }{M_i}}\vec{\bm u} + \boxed{\frac{\vec{\bm j}_i}{M_i}} \nonumber\\
            &\boxed{\vec{\bm N}_i = C_i \vec{\bm u} + \vec{\bm J}_i} \label{eq:molar_flux_decomposition}
          \end{align}
          衍生出了
        \item Species Molar Concentration
        \begin{equation}
          \boxed{C_i = \frac{\rho_i}{M_i}}  \label{eq:species_molar_concentration}
        \end{equation}
        \item Molar Diffusion Flux
        \begin{equation}
          \boxed{\vec{\bm J}_i = \frac{\vec{\bm j}_i}{M_i}}
        \end{equation}
        \item Molar Rate of Production
        \begin{equation}
          \boxed{R_i = \frac{r_i}{M_i}} \label{eq:molar_rate_of_production}
        \end{equation}
        注意因為是反應莫爾數，因此加總會跟總反應的係數和有關
        \begin{equation}
          \sum_i R_i \neq 0 (\text{Depends on Stoichiometry})
        \end{equation}
      \end{enumerate}
      但若使用(\ref{eq:molar_flux_decomposition})來計算總莫耳通量，會發現因為$\sum_i \vec{\bm j}_i = 0$
      \begin{equation}
        \sum_i \vec{\bm N}_i = \sum_i C_i \vec{\bm u}_i \label{eq:species_molar_flux_alternate}
      \end{equation}
      移項候可以用此表達$\vec{\bm j}_i$
      \begin{align}
        \vec{\bm N}_i &= C_i \vec{\bm u} + \vec{\bm J}_i \nonumber\\
        &= C_i \vec{\bm u}_i \nonumber\\
        \implies \vec{\bm J}_i &= C_i (\vec{\bm u}_i - \vec{\bm u})
      \end{align}
      此時若將$\vec{\bm u}$的定義:
      \begin{equation}
        \vec{\bm u} = \frac{\sum_i \rho_i \vec{\bm u}}{\sum_i \rho_i} = \frac{\sum_i C_i M_i \vec{\bm u}_i}{\sum_i C_i M_i}
      \end{equation}
      代入會發現
      \begin{align}
        \sum_i \vec{\bm J}_i &= \sum_i C_i (\vec{\bm u}_i - \vec{\bm u}) \nonumber\\
        &= \sum_i C_i \vec{\bm u}_i - \vec{\bm u} \sum_i C_i \nonumber\\
        &= \sum_i C_i \vec{\bm u}_i - \vec{\bm u} C \neq 0
      \end{align}
      也就是說:
      \begin{equation}
        \sum_i \vec{\bm J}_i \neq 0
      \end{equation}
      為了使得$\sum_i \vec{\bm J}_i = 0$，因此需要\fbox{重新定義一個平均流速}
      \item Molar Average Velocity, $\vec{\bm u}^\ast$
        因為$\sum_i \vec{\bm J}_i \neq 0$，因此定義新的$\vec{\bm u}^\ast$
        \begin{equation}
          \boxed{C \vec{\bm u}^\ast = \sum_i C_i \vec{\bm u}_i \implies \vec{\bm u}^\ast = \frac{\sum_i C_i \vec{\bm u}_i}{C} }
          \label{eq:molar_flux_decomposition_volume_average}
        \end{equation}
        能使得:
        \begin{equation}
          \sum_i \vec{\bm J}_i = 0
        \end{equation}
        且
        \begin{equation}
          \vec{\bm J}_i = C_i \left(\vec{\bm u}_i - \vec{\bm u}^\ast\right)
        \end{equation}
    \end{enumerate}
\item 將動量傳輸、能量傳輸、質量傳輸的通量做比較
\begin{align}
  \text{Momentum Flux: }& \overline{\overline {\bm \Phi}} = \underbrace{\quad\rho \vec{\bm u}\vec{\bm u}\quad}_{\text{慣性力}} 
  + \underbrace{\quad\quad\overline{\overline{\bm \pi}}\quad\quad}_{\text{周遭流體作用力}} \nonumber\\
  \text{Energy Flux: }& \overline{\bm e}  = \underbrace{\quad\rho\hat E \vec{\bm u} \quad}_{\text{對流}}
  + \underbrace{\quad\vec{\bm_q}\quad}_{\text{熱傳導}} 
  + \underbrace{\quad\quad\overline{\overline{\bm \pi}}\cdot \vec{\bm u}\quad\quad}_{\text{分子運動做功}} \nonumber\\
  \text{Mass Flux: }& n_i = \underbrace{\quad\rho_i \vec{\bm u}\quad}_{\text{對流}}
  + \underbrace{\quad\vec{\bm j}_i\quad}_{\text{分子擴散}} \label{eq:transport_flux_comparison}
\end{align}
\item 熱傳與質傳的比較筆記:
\renewcommand{\arraystretch}{1.5}
\begin{table}[H]
  \centering
  \begin{tabular}{|c|c|c|}
    \hline
    項目 & 熱傳 & 質傳 \\
    \hline
    Govering equation &
    $\alpha \frac{\partial^2 T}{\partial x^2}= \frac{\partial T}{\partial t}$ &
    $D_{AB} \frac{\partial^2 C_A}{\partial x^2} = \frac{\partial C_A}{\partial t}$ \\
    \hline
    Dimension less time &
    $t^\ast \equiv \textrm{Fo} \equiv \frac{\alpha t}{L^2}$ &
    $t^\ast \equiv \textrm{Fo}_m \equiv \frac{D_{AB} t}{L^2}$ \\
    \hline
    Dimension less temp/conc &
    $\theta^\ast \equiv \frac{T-T_\infty}{T_s - T_\infty}$ &
    $\gamma^\ast \equiv \frac{C_A-C_{A,\infty}}{C_{A,i} - C_{A,\infty}}$ \\
    \hline
    Similarity variable &
    $\eta = \frac{x}{2\sqrt{\alpha t}}$ &
    $\eta = \frac{x}{2\sqrt{D_{AB} t}}$ \\
    \hline
    Biot number &
    $\textrm{Bi} = \frac{hL}{k}$ &
    $\textrm{Bi}_m = \frac{k_c L}{D_{AB}}$ \\
     \hline
  \end{tabular}
  \caption{熱傳與質傳的比較}
\end{table}
\item 擴散係數的影響因子(依照重要程度排序)
\begin{enumerate}
  \item 介質(分子間距離):\\
  氣體的擴散係數最大，液體次之，固體最小\\
  因為氣體分子間距較大，分子運動較自由，擴散較快\\
  反之固體分子間距較小，分子運動受限，擴散較慢
  \item 溫度: 溫度越高，分子熱運動越劇烈，擴散係數越大
  \begin{equation}
    D_{AB} \propto \sqrt{T} \quad \text{(氣體)}
  \end{equation}
  \item 分子量(分子移動速率):
  \begin{equation}
    v \propto \sqrt{\frac{T}{M}} \implies D_{AB} \propto \frac{1}{\sqrt{M}}
  \end{equation}
  \item 壓力(分子數量):\\
   壓力越高，分子間距越小，擴散係數越小
\end{enumerate}
故分子的擴散可以根據平均自由徑$\lambda$分成兩種\\
平均自由徑$\lambda$是指分子兩次碰撞之間的平均距離
\begin{enumerate}
  \item Molecular Diffusion(分子擴散):\\
  當$\lambda$遠小於系統的特徵長度$d$時，分子擴散佔主導地位\\
  此時分子間頻繁碰撞，擴散行為可用Fick's Law描述
  \item Knudsen Diffusion(克努森擴散):\\
  當例如在孔洞中，$\lambda$遠大於系統的特徵長度$d$時，克努森擴散佔主導地位\\
  那碰撞就不在跟其他分子，而是跟孔洞壁面碰撞\\
  此時擴散行為可用Knudsen擴散方程式描述
  \begin{equation}
    D_{K} = \frac{2}{3} r_p \sqrt{\frac{8RT}{\pi M}}
  \end{equation}
  其中$r_p$是孔洞半徑
\end{enumerate}
\item 質傳的四大定律(經驗式):
\begin{itemize}
  \item Film Theory(薄膜定律)
    \begin{equation}
      k_c \propto D_{AB}
    \end{equation}
  \item Penetration Theory(穿透定律)
    \begin{equation}
      k_c \propto \sqrt{D_{AB}}
    \end{equation}
  \item Boundary Layer Theory(邊界層定律)
    \begin{equation}
      k_c \propto D_{AB}^{\frac{2}{3}}
    \end{equation}
  \item Two Films Theory(雙薄膜定律)
    \begin{equation}
      \frac{1}{K_y} = \frac{1}{k_{y}} + \frac{m'}{k_x}
    \end{equation}
    或
    \begin{equation}
      \frac{1}{K_x} = \frac{1}{k_x} + \frac{1}{ky\cdot m''}
    \end{equation}
\end{itemize}
\item 無因次群:
\begin{itemize}
  \item Schmidt Number:
  \begin{equation}
    \text{Sc} = \frac{\nu}{D_{AB}} = \frac{\mu}{\rho D_{AB}} = \frac{\text{Momentum Diffusivity}}{\text{Mass Diffusivity}}
  \end{equation}
  這裡也同熱對流有:
  \begin{equation}
    \text{Sc} = \left(\frac{\delta}{\delta_m}\right)^3 = \left(\frac{\text{Mass Boundary Layer}}{\text{Momentum Boundary Layer}}\right)^3
  \end{equation}
  的關係，但注意$\delta_m$是質量邊界層厚度，而不是熱對流的熱邊界層厚度
  \item Lewis Number:
  \begin{equation}
    \text{Le} = \frac{\alpha}{D_{AB}} = \frac{k}{\rho C_p D_{AB}} = \frac{\text{Thermal Diffusivity}}{\text{Mass Diffusivity}}
  \end{equation}
  \item Prandtl Number:
  \begin{equation}
    \text{Pr} = \frac{\nu}{\alpha} = \frac{C_p\mu}{k} = \frac{\text{Momentum Diffusivity}}{\text{Thermal Diffusivity}}
  \end{equation}
  \item Sherwood Number:\\
  對比熱對流，就是那邊的Nusselt Number，徑向的動量傳輸和質量傳輸比
  \begin{equation}
    \text{Sh} = \frac{k_cL}{D_{AB}} = \frac{\text{Mass Convection Rate}}{\text{Molecular Diffusion Rate}}
  \end{equation}
  若系統為圓柱或球:
  \begin{equation}
    \text{Sh} = \frac{k_cD}{D_{AB}}
  \end{equation}
  \item Peclet Number:\\
  在軸向上的質量傳輸比
  \begin{equation}
    \text{Pe}_M =\text{Re}\cdot\text{Sc} = \frac{\rho u D}{\mu} = \frac{\text{Mass Convection Rate}}{\text{Molecular Diffusion Rate}}
  \end{equation}
  質量跟著液體流動方向傳遞\\
  若$\text{Pe}_M$很大，代表軸向上的質量傳導可以被忽略
  \item Stanton Number:
  \begin{equation}
    \text{St}_M = \frac{\text{Sh}}{\text{Re}\cdot\text{Sc}} = \frac{k_c}{\nu}
  \end{equation}
\end{itemize}
\item 一些質量傳輸的經驗式:
\begin{itemize}
  \item 流體流過平板:
  \begin{equation}
    \text{Sh} = 0.664\text{Re}^{\frac{1}{2}}\text{Sc}^{\frac{1}{3}}
  \end{equation}
  \item 層流流過長直圓管，其壁面濃度均一(泡茶包)
  \begin{equation}
    \text{Sh} = 3.66 + f(\text{Re},\text{Sc},\frac{L}{D})
  \end{equation}
  \item 層流流過長直圓管，給予均勻質量通量(過濾器)
  \begin{equation}
    \text{Sh} = \frac{48}{11} + f(\text{Re},\text{Sc},\frac{L}{D})
  \end{equation}
  \item Plug Flow，流過長直圓管，給予均勻質量通量
  \begin{equation}
    \text{Sh} = 8 + f(\text{Re},\text{Sc},\frac{L}{D})
  \end{equation}
  \item 流過球面的流(泡澡球):
  \begin{equation}
    \text{Sh} = 2 + f(\text{Re},\text{Sc})
  \end{equation}
  \item Reynolds Analogy:(適用於$\text{Sc}=1$)
  \begin{equation}
    \frac{k_c}{u_\infty} = \text{St}_M = \frac{C_f}{2} = \frac{\text{Sh}}{\text{Re}\cdot\text{Sc}}
  \end{equation}
  P.S. 回顧熱對流那邊，如果Pr也等於1，則有Complete Reynolds Analogy:
  \begin{equation}
    St_H = St_M = \frac{C_f}{2} = \frac{\text{Nu}}{\text{Re}\cdot\text{Pr}}
  \end{equation}
  也就是說當Sc=Pr=1時，熱質傳行為相同
  \begin{equation}
    \text{Sc} = \text{Pr} = 1 \implies St_H = St_M = \frac{C_f}{2}
  \end{equation}
  \item Colburn Analogy:(適用於$0.6<\text{Sc}<2500$)\\
   (P.S. 熱對流那邊是$0.5<\text{Pr}<50$)\\
   可以說是Modified Reynolds Analogy
  \begin{equation}
    j_D = \text{St}_M\cdot \text{Sc}^{\frac{2}{3}} = \frac{C_f}{2}
  \end{equation}
  $j_D$是Colburn $j$ factor for mass transfer\\
  所以有Complete Colburn Analogy:
  \begin{equation}
    j_H =  j_D = \frac{C_f}{2},\quad 0.6<\text{Sc}<2500,\quad 0.5<\text{Pr}<50
  \end{equation}
  \item Prandtl Analogy:(適用於忽略Drag)
  \begin{equation}
    \text{St}_M = \frac{\frac{C_f}{2}}{1+5\sqrt{\frac{C_f}{2}}\left(\text{Sc}-1\right)} = \frac{k_c}{u_\infty}
  \end{equation}
  \item Von Karman Analogy:(適用於忽略Drag)
  \begin{equation}
    \text{St}_M = \frac{\frac{C_f}{2}}{1+5\sqrt{\frac{C_f}{2}\left\{
      \text{Sc}-1+\ln\left[1+\frac{5}{6}\left(\text{Sc}-1\right)\right]
    \right\}}} = \frac{k_c}{u_\infty}
  \end{equation}
\end{itemize}
\end{itemize}
\end{CJK*}
\end{document}